
\exo 

\begin{enumerate}
\item on donne $p(A) = 0,4$  \qquad ; \qquad $p(B)=0,7$ \qquad et \qquad $p(A\cap B)=0,2$

Calculer $p(A \cup B)$
\item on donne $p(R) = 0,6$  \qquad ; \qquad $p(S)=0,8$ \qquad et \qquad $p(R\cup S)=0,9$

Calculer $p(R \cap S)$
\item on donne $p(E) = 0,6$  \qquad ; \qquad $p(E \cap F)=0,5$ \qquad et \qquad $p(E\cup F)=0,7$

Calculer $p(F)$


\end{enumerate}

\medskip

\exo

Une enquête réalisée par un journal local révèle que 48\% de leurs abonnés lisent à chaque fois la page "Spectacles et Animations" ("A"), que 67\% ne manquent pas la page "Sports" ("S") et que 27\% lisent toujours ces deux pages avec le même intérêt.

Calculer la probabilité qu'un abonnée pris au hasard :
\begin{enumerate}%[label=\roman{enumi}}
\item Représenter la situation par un diagramme de Venn.
\item lise au moins l'une des deux pages ;
\item ne lise pas la page "Spectacles et Animations" ;
\item lise la page "Sport" et pas la page "Spectacles et Animations"
\end{enumerate}

\medskip


\exo 

En France il y a 16\% de gauchers, 30\% de personnes aux yeux bleus et 3\% de gauchers aux yeux bleus.

Quelle est la probabilité qu'un individu pris au hasard ne soit pas gaucher et n'ait pas les yeux bleus ?

\medskip

\exo

Dans un lycée, on a effectué un sondage sur les activités extra-scolaires :

10\% des élèves font à la fois du sport et jouent d'un instrument de musique. 

25\% font seulement du sport et 40\% ne pratiquent ni sport ni musique.

On interroge un élève au hasard à la sortie du lycée.
\begin{enumerate}
\item Représenter la situation par un diagramme de Venn ou un tableau.
\item Déterminer les probabilités des événements suivants :
		\begin{enumerate}
		\item L'élève ne fait pas de sport, mais joue d'un instrument.
		\item L'élève joue d'un instrument.
		\item L'élève pratique un sport ou joue d'un instrument.
		\end{enumerate}
\end{enumerate}

\exo

Les 240 adhérents d'une association sportive d'un lycée doivent choisir un seul sport parmi le basket-ball, le ping-pong et  la natation.

Le tableau ci-dessous donne la répartition des adhérents de cette association en tenant compte du régime du lycée (Demi-Pensionnaire ou externe) :

\begin{center}
\begin{tabularx}{0.8\linewidth}{|*{4}{>{\centering \arraybackslash \footnotesize}X|}c|}\hline
& Basket-ball&Ping-pong&Natation&Total\\ \hline 
Demi-pensionnaires&30&40&60&130\\ \hline 
Externes&36&32&42&110\\ \hline 
Total&66&72&102&210\\ \hline 
\end{tabularx}
\end{center}

\emph{Dans les questions suivantes, on donnera les résultats sous forme décimale, arrondis au centième.}

\medskip
 
On choisit au hasard un adhérent de l'association. On suppose que tous les adhérents ont la même probabilité d'être choisis.

\medskip
\begin{enumerate}
 
\item Calculer la probabilité de chacun des événements suivants :
 
$A$ : \og l'adhérent a choisi le Basket-ball \fg{} ; 
$B$ : \og l'adhérent est externe \fg.

\item Exprimer l'événement $A \cup B $ à l'aide d'une phrase, puis calculer sa probabilité. 

\item Quelle est la probabilité que l'adhérent ne soit ni externe ni joueur de basket-ball ?
\item On interroge un externe, quelle est la probabilité qu'il fasse du Ping-pong ?

\end{enumerate} 

\medskip



\exo 

La gérante d'un hôtel 3 étoiles du IX\up{e} arrondissement de la ville de Paris analyse les
réservations effectuées pour le premier week-end du mois de juin.

\begin{itemize}
\item[$\bullet~~$] Cent chambres sont réservées: 55 chambres simples et 45 chambres doubles.
\item[$\bullet~~$] Parmi les chambres simples réservées, 40\,\% l'ont été par des clients ayant entre
18 et 30 ans, 20\,\% par des clients de plus de 51 ans et le reste par des clients ayant
entre 31 et 50 ans.
\item[$\bullet~~$] Un tiers des chambres doubles a été réservé par des clients ayant entre 18 et
30 ans, un neuvième par des clients ayant entre 31 et 50 ans, le reste par des clients de
plus de 51 ans.
\end{itemize}
\setlength\parindent{0mm}


\begin{enumerate}
\item Compléter le tableau :
\begin{center}
\begin{tabularx}{0.8\linewidth}{|m{1.75cm}|*{3}{>{\centering \arraybackslash}X|}c|}\hline
&Entre 18 et 30 ans &Entre 31 et 50 ans &Plus de 51 ans& Total\\ \hline
Chambre simple&&&&\\ \hline
Chambre double&&&&\\ \hline
Total &&&&100 \rule[-8pt]{0pt}{10pt}\\ \hline
\end{tabularx}
\end{center}


\item La gérante de l'hôtel choisit au hasard une fiche de réservation d'une chambre de
son établissement pour le premier week-end de juin. On admet que chaque fiche
possède la même probabilité d'être choisie. On considère les évènements suivants:
\begin{description}
\item[ ] $A$ : \og La fiche choisie est celle d'un client ayant entre 18 et 30 ans \fg.
\item[ ] $B$ : \og La fiche choisie est celle d'un client ayant entre 31 et 50 ans \fg.
\item[ ] $C$ : \og La fiche choisie est celle d'un client ayant plus de 51 ans \fg.
\item[ ] $D$ : \og La fiche choisie correspond à la réservation d'une chambre double \fg.
\end{description}

	\begin{enumerate}
		\item Déterminer la probabilité des événements A, B, C et D.
		\item Définir par une phrase l'évènement $B  \cap D$ puis calculer sa probabilité.
		\item Calculer la probabilité de l'évènement \og La fiche choisie est celle d'un client ayant entre 31 et 50 ans ou ayant réservé une chambre double \fg.
		\item Décrire, à l'aide d'une phrase, l'évènement $\overline{A}$ puis calculer la probabilité de $\overline{A}$.
		
 	\end{enumerate}

\item  La gérante a choisi une fiche de réservation parmi celles des clients ayant
entre 31 et 50 ans.
	
On admet que chacune de ces fiches possède la même probabilité d'être choisie.
	
Calculer la probabilité que la fiche choisie corresponde à la réservation d'une
chambre double.
\end{enumerate}

\medskip

\exo

Une urne contient une boule bleue(B), une boule blanche(L) et une boule rouge(R).

On extrait au hasard une boule de l'urne, on note sa couleur et on la remet dans l'urne.

Puis on recommence encore deux fois.

On obtient ainsi un tirage successif, avec remise, de trois boules.

\begin{enumerate}
\item Faire un arbre, et en déduire le nombre de résultats possibles.
\item  En déduire la probabilité de chacun des événements suivants :

A : \og les trois boules sont rouges \fg ;

B : \og les trois boules sont de même couleur \fg{};

C : \og au moins une boule est rouge \fg{} ;

D : \og on a extrait trois boules de couleurs différentes \fg{}.

\end{enumerate}

\medskip

\exo

Nicolas répond à un exercice "Vrai ou Faux" totalement au hasard. 
Une seule réponse est exacte.

Cet exercice est composé de 4 questions.

\begin{enumerate}
\item Faire un arbre puis donner le nombre de possibilités.
\item Quelle est la probabilité que Nicolas ait tout juste ?
\item Quelle est la réponse qu'il obtiennent au moins une bonne réponse ?

\item Reprendre l'énoncé(répondre les deux questions précédentes) lorsque l'exercice est un QCM comportant 3 réponses possibles et où une seule réponse est exacte.
\end{enumerate}

\exo

Une usine produit des \emph{smartphones}. Les contrôles ont fait apparaître que 5\% des \emph{smartphones} ont un défaut à l'écran tactile et 3\% ont un défaut à la batterie, tandis que 1\% ont les deux défauts.

On choisit un \emph{smartphone} produit au hasard.

On note $E$ l'événement \og le mobile a un défaut à l'écran \fg{}  et $B$ l'événement \og le mobile a un défaut à la batterie \fg{} .

\begin{enumerate}
\item Donner $p(E)=\ldots\ldots$\qquad\qquad et $p(B)=\ldots\ldots$.
\item Quelle est la probabilité que le téléphone n'ait pas de défaut à l'écran.

\item Déterminer la probabilité que le téléphone ait au moins un défaut.


\end{enumerate}

\medskip

\exo

Les 240 adhérents d'une association sportive d'un lycée doivent choisir un seul sport parmi le basket-ball, le ping-pong et  la natation.

Le tableau ci-dessous donne la répartition des adhérents de cette association en tenant compte du régime du lycée (Demi-Pensionnaire ou externe) :

\begin{center}
\begin{tabularx}{0.8\linewidth}{|*{4}{>{\centering \arraybackslash \footnotesize}X|}c|}\hline
& Basket-ball&Ping-pong&Natation&Total\\ \hline 
Demi-pensionnaires&30&40&60&130\\ \hline 
Externes&36&32&42&110\\ \hline 
Total&66&72&102&210\\ \hline 
\end{tabularx}
\end{center}

\emph{Dans les questions suivantes, on donnera les résultats sous forme décimale, arrondis au centième.}

\medskip
 
On choisit au hasard un adhérent de l'association. On suppose que tous les adhérents ont la même probabilité d'être choisis.

\medskip
\begin{enumerate}
 
\item Calculer la probabilité de chacun des événements suivants :
 
$A$ : \og l'adhérent a choisi le Basket-ball \fg{} ; 

$B$ : \og l'adhérent est externe \fg.

\item Exprimer l'événement $A \cup B $ à l'aide d'une phrase, puis calculer sa probabilité. 

\item Quelle est la probabilité que l'adhérent ne soit ni externe ni joueur de basket-ball ?

\item On interroge un externe, quelle est la probabilité qu'il fasse du Ping-pong ?

\end{enumerate} 

\exo

Une urne contient une boule bleue(B), une boule blanche(L) et une boule rouge(R).

On extrait au hasard une boule de l'urne, on note sa couleur et on la remet dans l'urne.

Puis on recommence encore deux fois.

On obtient ainsi un tirage successif, avec remise, de trois boules.

\begin{enumerate}
\item Faire un arbre, et en déduire le nombre de résultats possibles.

\item  En déduire la probabilité de chacun des événements suivants :

A : \og les trois boules sont rouges \fg ;

B : \og les trois boules sont de même couleur \fg{};

C : \og au moins une boule est rouge \fg{} ;

D : \og on a extrait trois boules de couleurs différentes \fg{}.

\end{enumerate}

\exo
Sophie a mis des dragée dans une boîte, les unes contiennent une amande, les autres n'en contiennent pas.\\
On sait que :
\begin{enumerate}
\item[$\bullet$] 30\% des dragées contiennent une amande ;
\item[$\bullet$] 40\% des dragées avec amande sont bleues, les autres sont roses ;
\item[$\bullet$] 75\% des dragées sans amande sont bleues, les autres sont roses.
\end{enumerate}
Sophie choisit au hasard une dragée dans la boîte. On admet que toutes les dragées ont la même probabilité d'être choisies.\\
On considère les évènements :
\begin{description}
\item $A$ : \og la dragée choisie contient une amande \fg ;
\item $B$ : \og la dragée choisie est bleue \fg.
\end{description}
\begin{multicols}{2}

\begin{enumerate}
\item Recopier et compléter l'arbre des fréquences ci-contre.
\item Quelle est la proportion de dragées bleues contenant une amande dans la boîte ?
\item Décrire l'événement $A \cap B$ par une phrase. Montrer que sa probabilité est égale à 0,12.
\item Calculer la probabilité de l'évènement B.
\item Décrire l'évènement $A \cup B$ par une phrase, puis calculer sa probabilité.
\end{enumerate}


\hspace*{2cm}
\begin{psmatrix}[rowsep=.3cm, colsep=1.5cm]
 & & $B$\\
 & $A$ & \\
 & & $\overline{B}$\\
 ~~&&&\\
 &&$B$\\
 &$\overline{A}$&\\
 &&$\overline{B}$\\
 
 
 \psset{linewidth=1pt}
\ncline{4,1}{2,2}^{$30\%$}
\ncline{2,2}{1,3}^{$40\%$}
\ncline{2,2}{3,3}_{$\ldots$}
\ncline{4,1}{6,2}_{$\ldots$}
\ncline{6,2}{5,3}^{$75\%$}
\ncline{6,2}{7,3}_{$\ldots$}

\end{psmatrix}

\end {multicols}


\medskip
\exo
En France, la proportion de naissances de garçons est de 51,2\%. \og On s'intéresse aux familles de trois enfants sans jumeaux. \fg \\
On codera une famille par un triplet. Par exemple, pour la famille ($ F ; G ; F $), l'aîné est une fille, le cadet est un garçon et le benjamin est une fille.
\begin{enumerate}
\item Construire un arbre "pondéré" traduisant la situation.
\item On choisit au hasard une famille de trois enfants sans jumeaux.
\begin{enumerate}
\item Déterminer la probabilité que cette famille ait trois garçons.
\item Déterminer la probabilité que cette famille ait exactement deux garçons.
\item Déterminer la probabilité que cette famille ait au moins deux garçons.
\end{enumerate}
\end{enumerate}

\medskip
\exo \begin{multicols} {2}
Un jeu consiste à faire tourner une roue, attendre son arrêt, puis lancer un dé équilibré à six faces. La roue ci-contre est bien équilibré : la probabilité que la flèche indique un secteur est proportionnelle à l'angle au centre de ce secteur.\\
La règle du jeu est la suivante :
\begin{description}
\item[-] si la flèche indique le secteur $A$ et si le dé sort un "6" alors le joueur gagne le gros lot ;
\item[-] si la flèche indique le secteur $B$ et si le dé sort un nombre impair, le joueur gagne le lot moyen ;
\item[-] dans tous les autres cas le joueur perd.
\end{description}

\hspace*{3cm}
\begin{pspicture} (-2,-2) (2,2)
\pswedge[fillstyle=solid,fillcolor=gray](0,0){2}{0}{90}
\definecolor{nuance}{gray}{0.6 }
\pswedge[fillstyle=solid,fillcolor=nuance](0,0){2}{90}{210}
\pswedge[fillstyle=solid,fillcolor=lightgray](0,0){2}{210}{360}
\SpecialCoor
\psline[linewidth=0.1cm]{->}(1;-120)(1;60)
\rput(1.3;120){$B$}
\rput(1.3;50){$A$}
\rput(1.3;-70){$C$}
\psarc{-}(0,0){0.7}{90}{210}
\rput(1;150){$120^\circ$}
\end{pspicture}

\end{multicols}
\begin{enumerate}
\item Déterminer, pour chaque secteur ( $A$, $B$ ou $C$) la probabilité que la flèche indique ce secteur.
\item Faire un arbre pondéré et répondre aux questions suivantes : 
\begin{enumerate}
\item Quelle est la probabilité de gagner le gros lot ?
\item Quelle est la probabilité de gagner un lot ?
\item Quelle est la probabilité de perdre ?
\end{enumerate}
\end{enumerate}

\exo
Au rayon \og image et son \fg d'un grand magasin, un téléviseur et un lecteur de DVD sont en promotion pendant une semaine. Une personne se présente :
\begin{itemize}
\item la probabilité qu'elle achète le téléviseur est 3/5;
\item la probabilité qu'elle achète le lecteur de DVD si elle achète le téléviseur est 7/10;
\item la probabilité qu'elle achète le lecteur de DVD si elle n'achète pas le téléviseur est 1/10.
\end{itemize} 
On désigne par T l'événement : \og la personne achète le téléviseur \fg et par L l'événement : \og la personne achète le lecteur de DVD \fg .

\begin{enumerate}
\item Traduire les données de l'énoncé à l'aide d'un arbre pondéré.
\item Déterminer les probabilités des évènements suivants (les résultats seront donnés sous forme de fractions) :
\begin{enumerate}
\item \og la personne achète les deux appareils \fg .
\item \og la personne achète le lecteur de DVD \fg .
\item \og la personne n'achète aucun des deux appareils \fg . 
\end{enumerate}
\item On suppose que la personne achète le lecteur de DVD, quelle est la probabilité qu'elle achète aussi le téléviseur.
\end{enumerate}

\exo : \\
Un malade souffrant d'angine va consulter son médecin. L'agent infectieux a 4 chances sur 5 d'être un streptocoque. Le médecin décide de faire des analyses en laboratoire. Les techniques de laboratoire comportent des risques d'erreur :
\begin{description}
\item[$\bullet$] Le streptocoque, lorsqu'il est présent, a 1 chance sur 5 de ne pas être décelé.
\item[$\bullet$] Le streptocoque, lorsqu'il est absent, a 1 chance sur 10 d'être décelé par erreur.

\end{description}

On note S  l'événement : \og Le streptocoque est présent \fg et  D l'événement  \og Le streptocoque est décelé \fg .

{\bf PARTIE A}

Dans cette partie, les résultats seront donnés sous forme de fractions irréductibles.

Le médecin fait procéder à une première analyse.
\begin{enumerate}
\item Traduire les données à l'aide d'un arbre pondéré.
\item Calculer la probabilité de l'événement : \og Le Streptocoque est présent et est décelé \fg .
\item Montrer que la probabilité P(D) que le streptocoque soit décelé est égale à $\dfrac{33}{50}$.
\item Le streptocoque est décelé. Quelle est la probabilité pour qu'il soit présent ?
\end{enumerate}

{\bf PARTIE B}

Pour confirmer son diagnostic, le médecin fait analyser quatre autres prélèvements faits sur ce patient.

(Les quatre tests sont réalisés dans les mêmes conditions et sont indépendants).

Quelle est la probabilité pour que le streptocoque soit décelé dans exactement trois tests parmi les quatre ?

